% Part: first-order-logic
% Chapter: tableaux
% Section: soundness.tex

\documentclass[../../../include/open-logic-section]{subfiles}

\begin{document}

\iftag{FOL}
      {\olfileid[pt]{fol}{tab}{sou}}
      {\olfileid[pt]{pl}{tab}{sou}}
      
\olsection{Correção}

\begin{explain}
!!^a{derivation} sistema, como o de tableaux, é \emph{correto}
se não pode !!{derive} coisas que de fato não valem. A correção é,
portanto, uma espécie de propriedade de segurança garantida dos sistemas de !!{derivation}.
Dependendo da propriedade da teoria da prova em questão, gostaríamos
de saber, por exemplo, que
\begin{enumerate}
\item todo $!A$ !!{derivable} é \iftag{FOL}{válido}{uma tautologia};
\item se !!a{sentence} é !!{derivable} a partir de algumas outras, ela também é uma
  consequência delas;
\item se um conjunto de !!{sentence}s é inconsistente, ele é insatisfatível.
\end{enumerate}
Essas são propriedades importantes de !!a{derivation} sistema. Se alguma delas não
vale, o sistema de !!{derivation} é deficiente---ele !!{derive} coisas demais.
Consequentemente, estabelecer a correção de !!a{derivation} sistema é de
máxima importância.

Como todas essas propriedades da teoria da prova são definidas por meio de
!!{tableau}s fechados de um tipo ou de outro, demonstrar (1)--(3) acima exige
demonstrar algo sobre as propriedades semânticas dos
!!{tableau}s fechados. Primeiro definiremos o que significa !!a{signed
  formula} ser satisfeita em uma estrutura e, em seguida, mostraremos que, se um
!!{tableau} é fechado, nenhuma estrutura satisfaz todas as suas suposições.
(1)--(3) seguem então como corolários desse resultado.
\end{explain}

\begin{defn}
\iftag{FOL}{!!^a{structure}}{!!^a{valuation}}~$\iftag{FOL}{\Struct{M}}{\pAssign{v}}$
\emph{satisfaz} !!a{signed formula} $\sFmla{\True}{!A}$ sse
$\iftag{FOL}{\Sat{M}}{\pSat{v}}{!A}$, e satisfaz
$\sFmla{\False}{!A}$ sse
$\iftag{FOL}{\Sat/{M}}{\pSat/{v}}{!A}$. $\iftag{FOL}{\Struct{M}}{\pAssign{v}}$
satisfaz um conjunto de !!{signed formula}s~$\Gamma$ sse satisfaz toda
$\sFmla{S}{!A} \in \Gamma$. $\Gamma$~é \emph{satisfatível} se existe
\iftag{FOL}{!!a{structure}}{!!a{valuation}} que o satisfaz, e é
\emph{insatisfatível} caso contrário.
\end{defn}

\begin{thm}[Correção]
  \ollabel{thm:tableau-soundness}
  Se $\Gamma$ tem um !!{tableau} fechado, $\Gamma$ é insatisfatível.
\end{thm}

\begin{proof}
Chamemos um ramo de !!a{tableau} de satisfatível sse o conjunto de
!!{signed formula}s nele é satisfatível, e chamemos !!a{tableau} de
satisfatível se ele contém pelo menos um ramo satisfatível.

Mostramos o seguinte: estender um !!{tableau} satisfatível por uma das
regras de inferência sempre resulta em um !!{tableau} satisfatível.
Isso demonstrará o teorema: qualquer !!{tableau} fechado resulta da
aplicação de regras de inferência ao !!{tableau} que consiste apenas em
suposições de~$\Gamma$. Assim, se $\Gamma$ fosse satisfatível, qualquer
!!{tableau} para ele seria satisfatível. Um !!{tableau} fechado,
porém, claramente não é satisfatível: todo ramo contém tanto
$\sFmla{\True}{!A}$ quanto $\sFmla{\False}{!A}$, e nenhuma estrutura pode
ao mesmo tempo satisfazer e não satisfazer~$!A$.

Suponha que tenhamos um !!{tableau} satisfatível, O.s.\@, !!a{tableau} com pelo
menos um ramo satisfatível. Aplicar uma regra de inferência acrescenta
!!{signed formula}s a um ramo ou divide um ramo em dois. Se o
!!{tableau} tem um ramo satisfatível que não é estendido pela aplicação
da regra em questão, ele permanece um ramo satisfatível no
!!{tableau} estendido; portanto, o tableau estendido é satisfatível. Assim,
só precisamos considerar o caso em que uma regra é aplicada a um
ramo satisfatível.

Seja $\Gamma$ o conjunto de !!{signed formula}s nesse ramo, e
seja $\sFmla{S}{!A} \in \Gamma$ a !!{signed formula} à qual a
regra é aplicada. Se a regra não resulta na divisão do ramo, precisamos
mostrar que o ramo estendido, O.s.\@, $\Gamma$ juntamente com
as conclusões da regra, ainda é satisfatível. Se a regra resulta
na divisão do ramo, precisamos mostrar que pelo menos um dos dois
ramos resultantes é satisfatível.
  
Primeiro, consideramos as inferências possíveis que não resultam na divisão de um ramo.
\begin{enumerate}
\item O ramo é expandido aplicando-se $\TRule{\True}{\lnot}$ a
  $\sFmla{\True}{\lnot !B} \in \Gamma$. O ramo estendido
  contém então as !!{signed formula}s $\Gamma \cup
  \{\sFmla{\False}{!B}\}$. Suponha
  $\iftag{FOL}{\Sat{M}}{\pSat{v}}{\Gamma}$. Em particular,
  $\iftag{FOL}{\Sat{M}}{\pSat{v}}{\lnot !B}$. Portanto,
  $\iftag{FOL}{\Sat/{M}}{\pSat/{v}}{!B}$, O.s.\@,
  $\iftag{FOL}{\Struct{M}}{\pAssign{v}}$ satisfaz
  $\sFmla{\False}{!B}$.
\item O ramo é expandido aplicando-se $\TRule{\False}{\lnot}$ a
  $\sFmla{\False}{\lnot !B} \in \Gamma$: exercício.
\item O ramo é expandido aplicando-se $\TRule{\True}{\land}$ a
  $\sFmla{\True}{!B \land !C} \in \Gamma$, o que resulta em duas novas
  !!{signed formula}s no ramo: $\sFmla{\True}{!B}$ e
  $\sFmla{\True}{!C}$. Suponha
  $\iftag{FOL}{\Sat{M}}{\pSat{v}}{\Gamma}$, em particular
  $\iftag{FOL}{\Sat{M}}{\pSat{v}}{!B \land !C}$. Então,
  $\iftag{FOL}{\Sat{M}}{\pSat{v}}{!B}$ e
  $\iftag{FOL}{\Sat{M}}{\pSat{v}}{!C}$. Isso significa que
  $\iftag{FOL}{\Struct{M}}{\pAssign{v}}$ satisfaz tanto
  $\sFmla{\True}{!B}$ quanto $\sFmla{\True}{!C}$.
\item O ramo é expandido aplicando-se $\TRule{\False}{\lor}$ a
  $\sFmla{\False}{!B \lor !C} \in \Gamma$: exercício.
\item O ramo é expandido aplicando-se $\TRule{\False}{\lif}$ a
  $\sFmla{\False}{!B \lif !C} \in \Gamma$: isso resulta em duas novas
  !!{signed formula}s no ramo: $\sFmla{\True}{!B}$ e
  $\sFmla{\False}{!C}$. Suponha
  $\iftag{FOL}{\Sat{M}}{\pSat{v}}{\Gamma}$, em particular
  $\iftag{FOL}{\Sat/{M}}{\pSat/{v}}{!B \lif !C}$. Então,
  $\iftag{FOL}{\Sat{M}}{\pSat{v}}{!B}$ e
  $\iftag{FOL}{\Sat/{M}}{\pSat/{v}}{!C}$. Isso significa que
  $\iftag{FOL}{\Struct{M}}{\pAssign{v}}$ satisfaz tanto
  $\sFmla{\True}{!B}$ quanto $\sFmla{\False}{!C}$.

\iftag{FOL}{%
\item O ramo é expandido aplicando-se $\TRule{\True}{\lforall}$ a
  $\sFmla{\True}{\lforall[x][!B(x)]} \in \Gamma$: isso resulta em uma
  nova !!{signed formula}~$\sFmla{\True}{!A(t)}$ no ramo.
  Suponha $\Sat{M}{\Gamma}$, em particular,
  $\Sat{M}{\lforall[x][!A(x)]}$. Por
  \olref[syn][sem]{prop:quant-terms}, $\Sat{M}{!A(t)}$. Consequentemente,
  $\Struct{M}$ satisfaz $\sFmla{\True}{!A(t)}$.

\item O ramo é expandido aplicando-se $\TRule{\False}{\lforall}$ a
  $\sFmla{\False}{\lforall[x][!B(x)]} \in \Gamma$: isso resulta em uma
  nova !!{signed formula}~$\sFmla{\False}{!A(a)}$, em que $a$ é
  !!a{constant} que não ocorre em~$\Gamma$. Como $\Gamma$ é
  satisfatível, existe uma $\Struct{M}$ tal que $\Sat{M}{\Gamma}$, em
  particular $\Sat/{M}{\lforall[x][!B(x)]}$. Precisamos mostrar que
  $\Gamma \cup \{\sFmla{\False}{!A(a)}\}$ é satisfatível. Para isso,
  definimos uma $\Struct{M'}$ adequada da seguinte maneira.

  Por \olref[syn][ass]{prop:sat-quant}, $\Sat/{M}{\lforall[x][!B(x)]}$
  sse, para algum $s$, $\Sat/{M}{!B(x)}[s]$. Agora, seja $\Struct{M'}$
  igual a $\Struct{M}$, exceto que $\Assign{a}{M'} = s(x)$. Por
  \olref[syn][ext]{cor:extensionality-sent}, para qualquer
  $\sFmla{\True}{!C} \in \Gamma$, $\Sat{M'}{!C}$, e, para qualquer
  $\sFmla{\False}{!C} \in \Gamma$, $\Sat/{M'}{!C}$, pois $a$ não
  ocorre em~$\Gamma$.

  Por \olref[syn][ext]{prop:extensionality}, $\Sat/{M'}{!A(x)}[s]$. Por
  \olref[syn][ext]{prop:ext-formulas}, $\Sat/{M'}{!A(a)}[s]$. Como
  $!A(a)$ é !!a{sentence}, por
  \olref[syn][ass]{prop:sentence-sat-true}, $\Sat/{M'}{!A(a)}$, O.s.\@,
  $\Struct{M'}$ satisfaz $\sFmla{\False}{!A(a)}$.
\item O ramo é expandido aplicando-se $\TRule{\True}{\lexists}$ a
  $\sFmla{\True}{\lexists[x][!B(x)]} \in \Gamma$: exercício.
\item O ramo é expandido aplicando-se $\TRule{\False}{\lexists}$ a
  $\sFmla{\False}{\lexists[x][!B(x)]} \in \Gamma$: exercício.}{}
\end{enumerate}
Consideremos agora as inferências possíveis que resultam na divisão de um ramo.
\begin{enumerate}
\item O ramo é expandido aplicando-se $\TRule{\False}{\land}$ a
  $\sFmla{\False}{!B \land !C} \in \Gamma$, o que resulta em dois
  ramos: um à esquerda, que continua por $\sFmla{\False}{!B}$, e outro
  à direita, por $\sFmla{\False}{!C}$. Suponha
  $\iftag{FOL}{\Sat{M}}{\pSat{v}}{\Gamma}$, em particular
  $\iftag{FOL}{\Sat/{M}}{\pSat/{v}}{!B \land !C}$. Então,
  $\iftag{FOL}{\Sat/{M}}{\pSat/{v}}{!B}$ ou
  $\iftag{FOL}{\Sat/{M}}{\pSat/{v}}{!C}$. No primeiro caso,
  $\iftag{FOL}{\Struct{M}}{\pAssign{v}}$ satisfaz
  $\sFmla{\False}{!B}$, O.s.\@, $\iftag{FOL}{\Struct{M}}{\pAssign{v}}$
  satisfaz as fórmulas do ramo à esquerda. No segundo,
  $\iftag{FOL}{\Struct{M}}{\pAssign{v}}$ satisfaz
  $\sFmla{\False}{!C}$, O.s.\@, $\iftag{FOL}{\Struct{M}}{\pAssign{v}}$
  satisfaz as fórmulas do ramo à direita.
\item O ramo é expandido aplicando-se $\TRule{\True}{\lor}$ a
  $\sFmla{\True}{!B \lor !C} \in \Gamma$: exercício.
\item O ramo é expandido aplicando-se $\TRule{\True}{\lif}$ a
  $\sFmla{\True}{!B \lif !C} \in \Gamma$: exercício.
\item O ramo é expandido por \Cut: isso resulta em dois ramos,
  um contendo $\sFmla{\True}{!B}$ e o outro contendo
  $\sFmla{\False}{!B}$. Como $\iftag{FOL}{\Sat{M}}{\pSat{v}}{\Gamma}$
  e ou $\iftag{FOL}{\Sat{M}}{\pSat{v}}{!B}$ ou
  $\iftag{FOL}{\Sat/{M}}{\pSat/{v}}{!B}$,
  $\iftag{FOL}{\Struct{M}}{\pAssign{v}}$ satisfaz o ramo à esquerda ou
  o ramo à direita.
\end{enumerate}
\end{proof}

\tagprob{FOL}
\begin{prob}
Complete a demonstração de \olref[fol][tab][sou]{thm:tableau-soundness}.
\end{prob}
\tagendprob

\tagprob{notFOL}
\begin{prob}
Complete a demonstração de \olref[pl][tab][sou]{thm:tableau-soundness}.
\end{prob}
\tagendprob

\begin{cor}
\ollabel{cor:weak-soundness}
Se $\Proves !A$, então $!A$ é \iftag{FOL}{válida}{uma tautologia}.
\end{cor}

\begin{cor}
\ollabel{cor:entailment-soundness}
Se $\Gamma \Proves !A$, então $\Gamma \Entails !A$.
\end{cor}

\begin{proof}
  Se $\Gamma \Proves !A$, então, para alguns $!B_1$, \dots, $!B_n \in
  \Gamma$, $\{\sFmla{\False}{!A}, \sFmla{\True}{!B_1}, \dots,
  \sFmla{\True}{!B_n}\}$ tem um !!{tableau} fechado. Por
  \olref{thm:tableau-soundness}, toda
  \iftag{FOL}{!!{structure}}{!!{valuation}}~$\iftag{FOL}{\Struct{M}}{\pAssign{v}}$
  torna algum $!B_i$ falso ou torna $!A$ verdadeiro. Portanto, se
  $\iftag{FOL}{\Sat{M}}{\pSat{v}}{\Gamma}$, então também
  $\iftag{FOL}{\Sat{M}}{\pSat{v}}{!A}$.
\end{proof}

\begin{cor}
\ollabel{cor:consistency-soundness}
Se $\Gamma$ é satisfatível, então é consistente.
\end{cor}

\begin{proof}
Demonstramos a contrapositiva. Suponha que $\Gamma$ não seja consistente.
Então existem $!B_1$, \dots, $!B_n \in \Gamma$ e um
!!{tableau} fechado para $\{\sFmla{\True}{!B_1}, \dots, \sFmla{\True}{!B_n}\}$. Por
\olref{thm:tableau-soundness}, não existe
$\iftag{FOL}{\Struct{M}}{\pAssign{v}}$ tal que
$\iftag{FOL}{\Sat{M}}{\pSat{v}}{!B_i}$ para todo $i=1$, \dots,~$n$. Mas,
nesse caso, $\Gamma$ não é satisfatível.
\end{proof}

\end{document}
